\section{Proof Details}
\label{app:proofs}

This appendix records the algebraic interfaces used by both final theorems.
Each component described here is instantiated in the Lean proof chain; none is
left as a premise of the six public projected entry points.

\subsection{Common population identities}

Write $d:=b-\beta^*$ and
\[
  \bar\Sigma:=\frac1m\sum_{e=1}^m\Sigma_e,
  \qquad
  Q(w):=A(w)+\frac{\bar\Sigma}{1+\gamma m}.
\]
Because
$\sum_e(w_e-a_\gamma)=1/(1+\gamma m)$, the corrected risk expansion gives
\begin{align}
  \mathcal L(b,w)
  &=\frac{\sigma_Y^2+2v_{\rm SEM}^\top d}{1+\gamma m}
    +d^\top Q(w)d,
  \label{eq:appendix-L-expansion}\\
  \nabla_b\mathcal L(b,w)
  &=\frac{2v_{\rm SEM}}{1+\gamma m}+2Q(w)d.
  \label{eq:appendix-L-gradient}
\end{align}
Hence
\begin{equation}
  \mathcal L(b,w)
  =\frac12d^\top\nabla_b\mathcal L(b,w)
   +\frac{\sigma_Y^2+v_{\rm SEM}^\top d}{1+\gamma m}.
  \label{eq:appendix-loss-gradient-identity}
\end{equation}
Since the penalty is independent of $b$,
\begin{equation}
  \mathcal L_\mu(b,w)
  =\frac12d^\top\nabla_b\mathcal L_\mu(b,w)
   +\frac{\sigma_Y^2+v_{\rm SEM}^\top d}{1+\gamma m}
   -\mu\|w\|_2^2.
  \label{eq:appendix-penalized-identity}
\end{equation}
These are the scalar identities formalized for the comparator reduction in
this note.  Their role is that the common term is identical at $w^0$ and
$w_t$ and therefore cancels in the subtraction.  They should not be read as
a renumbering of Eqs.~(64)--(67) in the official appendix, whose displayed
statements serve different purposes.

\subsection{Curvature and comparator signs}

At $w^0$, $A(w^0)\succeq\lambda I$ and
$\bar\Sigma\succeq0$.  Consequently
\begin{align*}
  d^\top\nabla_b\mathcal L(b,w^0)
  &\ge2\lambda\|d\|_2^2
      -\frac{2\|v_{\rm SEM}\|_2}{1+\gamma m}\|d\|_2\\
  &\ge\lambda\|d\|_2^2
      -\frac{\|v_{\rm SEM}\|_2^2}
             {\lambda(1+\gamma m)^2}.
\end{align*}
The Lean proof derives this from the actual symmetric second-moment matrices,
the Hermitian eigenvalue lower bound, and a finite-coordinate Young
inequality.

For the unpenalized comparator, subtraction of
\eqref{eq:appendix-loss-gradient-identity} gives
\begin{equation}
  d_t^\top\nabla_b\mathcal L(b_t,w_t)
  =d_t^\top\nabla_b\mathcal L(b_t,w^0)
   -2\operatorname{Reg}_t^0.
  \label{eq:appendix-unpenalized-comparator}
\end{equation}
For the penalized comparator,
\eqref{eq:appendix-penalized-identity} gives
\begin{align}
  d_t^\top\nabla_b\mathcal L_\mu(b_t,w_t)
  ={}&d_t^\top\nabla_b\mathcal L_\mu(b_t,w^0)
      -2\operatorname{Reg}_{t,\mu}^0\notag\\
     &-2\mu\bigl(\|w^0\|_2^2-\|w_t\|_2^2\bigr).
  \label{eq:appendix-penalized-comparator}
\end{align}
The regret sign is thus comparator minus current.  Only a cumulative upper
bound is used.

\subsection{Formalized entropy mirror-ascent chain}

Let $\phi(x)=x\log x$ with the standard zero convention, and let $D_\Phi$ be
the coordinate Bregman divergence generated by
$\sum_i\phi(w_i)$.  For positive simplex vectors it equals the finite natural-
logarithm relative entropy.  The formalized finite Pinsker theorem has the
normalization
\begin{equation}
  \frac12\left(\sum_i|p_i-q_i|\right)^2\le D_\Phi(p,q).
  \label{eq:appendix-pinsker}
\end{equation}
There is no factor $1/2$ in the definition of the $\ell_1$ distance.

The normalized exponentiated-gradient ascent update is strictly positive and
preserves the simplex.  For a coordinate bound $|g_i|\le G$, finite H\"older
gives
\[
  |\langle g,x\rangle|\le G\sum_i|x_i|.
\]
This is the finite-coordinate specialization of the standard entropy
mirror-ascent analysis; see, for example, Shalev-Shwartz~\cite{online-learning}.
The project proves the specialization directly rather than importing its
conclusion as an assumption.
Combining the exact three-point identity, \eqref{eq:appendix-pinsker}, and
\[
  -\frac12a^2+ba\le\frac12b^2
\]
proves the one-step inequality
\begin{equation}
  \eta_w\langle g,u-w\rangle
  \le
  D_\Phi(u,w)-D_\Phi(u,w^+)
  +\frac{\eta_w^2G^2}{2}.
  \label{eq:appendix-mirror-one-step}
\end{equation}
The comparator $u$ need only belong to the simplex; it need not be strictly
positive.  A separate tangent inequality proves
$D_\Phi(u,w^+)\ge0$ for such a boundary comparator because $w^+$ is positive.

Summing \eqref{eq:appendix-mirror-one-step} for the actual trajectory,
using the uniform initialization, and applying the initial-potential bound
\[
  D_\Phi(u,m^{-1}\mathbf1)\le\log m
\]
yields the pathwise estimate
\begin{equation}
  \sum_{t=1}^T\langle g_t,u-w_t\rangle
  \le
  \frac{\log m}{\eta_w}
  +\frac{\eta_w}{2}\sum_{t=1}^T G_t^2.
  \label{eq:appendix-entropy-regret}
\end{equation}
The predictable conditional-expectation bridge proves, coordinatewise and
then after finite summation, that
\[
  \E[\langle\widehat g_t,d_t\rangle\mid\mathcal F_{t-1}]
  =\langle F_t,d_t\rangle\quad\text{a.e.},
\]
and hence that the corresponding unconditional integrals are equal.  It does
not identify the two raw pairings a.e.  This is the formal route to
\eqref{eq:penalized-static-regret} and
\eqref{eq:unpenalized-static-regret}; no independence between the two oracle
vectors enters.

\subsection{Projection and regret-to-error reduction}

For a nonempty closed convex $C$, the constructed nearest-point map satisfies
the strong inequality
\begin{equation}
  \|P_C(x)-z\|_2^2+\|x-P_C(x)\|_2^2
  \le\|x-z\|_2^2,
  \qquad z\in C.
  \label{eq:appendix-pythagorean}
\end{equation}
Dropping the nonnegative middle term gives the Fej\'er inequality used in the
recursion.  Measurability of $P_C$ and the post-observation oracles proves
trajectory predictability by induction.

The common scalar reduction can be stated as follows.
\begin{proposition}[Comparator regret implies averaged squared error]
\label{prop:appendix-regret-to-error}
Suppose
\begin{align*}
  \E\|b_{t+1}-\beta^*\|_2^2
  \le{}&\E\|b_t-\beta^*\|_2^2
  -2\eta_b\E[(b_t-\beta^*)^\top F_t]
  +\eta_b^2G_b^2,\\
  (b_t-\beta^*)^\top F_t
  \ge{}&\lambda\|b_t-\beta^*\|_2^2
  -\varepsilon-\delta-2r_t,
\end{align*}
and $\sum_{t=1}^T\E[r_t]\le\mathcal R_T$.  Then
\begin{equation}
  D_T
  \le
  \frac{\|b_{\rm Init}-\beta^*\|_2^2}{2\lambda\eta_bT}
  +\frac{\varepsilon+\delta}{\lambda}
  +\frac{2\mathcal R_T}{\lambda T}
  +\frac{\eta_bG_b^2}{2\lambda}.
  \label{eq:appendix-abstract-bound}
\end{equation}
The same right-hand side bounds
$\E\|\bar b_T-\beta^*\|_2^2$.
\end{proposition}

\begin{proof}
Insert the direction inequality into the one-step recursion, sum from $1$ to
$T$, discard the nonnegative terminal squared distance, and divide by
$2\lambda\eta_bT$.  Finite-dimensional Jensen gives the averaged-iterate
statement.
\end{proof}

For the unpenalized result take
$\varepsilon=\varepsilon_{\rm id}$, $\delta=0$,
$r_t=\operatorname{Reg}_t^0$, and use
\eqref{eq:unpenalized-static-regret}.  For the penalized result take
$\delta=2\mu$, $r_t=\operatorname{Reg}_{t,\mu}^0$, and use
\eqref{eq:penalized-static-regret}.  These substitutions give the exact final
constants, including the single $2\mu/\lambda$ penalty term.

\subsection{Rate normalization and proof boundary}

For positive integer $T$,
\[
  \frac1{(1/\sqrt T)T}=\frac1{\sqrt T}.
\]
Substitution of $\eta_b=\eta_w=1/\sqrt T$ in the two finite-time bounds
therefore produces exactly $C_0/\sqrt T$ and $C_\mu/\sqrt T$.  The Lean source
proves this normalization as a real-number identity rather than invoking
asymptotic notation.

The proof audit covers the corrected population expansion, comparator signs,
finite Pinsker normalization, exponentiated-gradient regret, actual Euclidean
projection, conditional predictable pairings, oracle moment and integrability
bridges, stochastic recursion, spectral curvature, telescoping, Jensen, and
rate algebra.  It does not assert high-probability concentration,
last-iterate convergence, construction of a canonical fresh sample stream, or
verification of the official empirical exact-maximization algorithm.
